\documentclass[11pt]{article}
\usepackage[margin=0.82in]{geometry}
\usepackage{amsmath,amssymb,mathtools}
\usepackage{booktabs,tabularx}
\usepackage[dvipsnames]{xcolor}
\usepackage[colorlinks=true,linkcolor=NavyBlue,urlcolor=BrickRed,hypertexnames=false]{hyperref}
\usepackage{microtype,fancyhdr,enumitem}
\definecolor{ink}{HTML}{252E43}
\definecolor{accent}{HTML}{4B63B7}
\definecolor{pale}{HTML}{EFF2FF}
\definecolor{passgreen}{HTML}{16794A}
\color{ink}
\pagestyle{fancy}\fancyhf{}
\lhead{Symmetric groups}\rhead{Nonuniform permutation measures}\cfoot{\thepage}
\setlength{\headheight}{14pt}
\newcommand{\Sfun}{\mathcal S}
\newcommand{\Tr}{\operatorname{Tr}}
\newcommand{\PASS}{\textcolor{passgreen}{\textbf{PASS}}}
\newcommand{\summarybox}[1]{\begin{center}\fcolorbox{accent}{pale}{\parbox{0.91\linewidth}{#1}}\end{center}}
\title{\textbf{Nonuniform $\Lambda$-Distributions on Symmetric Groups}\\[3pt]
\large Cycle transforms, shuffles, and Fourier operators}
\author{Proof and computational verification note}
\date{17 July 2026}

\begin{document}
\maketitle
\summarybox{\textbf{Outcome.}  Explicit permutation matrices for
$S_3,S_4,S_5$ were compared with the cycle-count formula under six families
of nonuniform measures.  A separate $S_3$ calculation decomposes every tested
tensor-symmetric character into all three irreducibles and verifies both
central and noncentral Fourier operators.  All \textbf{267} checks pass; the
maximum absolute error is $4.45\times10^{-15}$.}

\section{Permutation matrices and the cycle transform}
For the defining permutation representation $V=\mathbb C^n$, a $d$-cycle
contributes all $d$th roots of unity.  Consequently
\begin{equation}
\det(I-tP_\pi)=\prod_{d\geq1}(1-t^d)^{C_d(\pi)}.            \tag{1}
\end{equation}
With
\begin{equation}
Q_d(\mathbf t)=\prod_a(1-t_a^d)^{-1},                      \tag{2}
\end{equation}
direct substitution in the determinant definition proves
\begin{equation}
\Sfun_{\nu,n}(\mathbf t)=
\mathbb E_\nu\!\left[\prod_{d\geq1}Q_d(\mathbf t)^{C_d(\pi)}\right]. \tag{3}
\end{equation}
This identity is valid for every probability measure, whether or not it is a
class measure.  It is the joint cycle-count probability-generating function
evaluated at the transformed variables $Q_d$.

For coefficient extraction, each factor is expanded using
\begin{equation}
(1-t^d)^{-c}=\sum_{j\geq0}\binom{c+j-1}{j}t^{dj}.           \tag{4}
\end{equation}
The verification uses integer dynamic programming for (4), separately for
every permutation.  This route does not diagonalize a matrix.

\section{Measures tested}
All probabilities are constructed on the fully enumerated group:
\begin{enumerate}[leftmargin=*,itemsep=3pt]
\item the lazy random-transposition walk with
$\kappa=0.35\delta_e+0.65\operatorname{Unif}(C_2)$, after three steps;
\item the lazy random $3$-cycle walk with laziness $0.4$, after two steps;
\item the lazy adjacent-transposition walk with laziness $0.25$, after four steps;
\item the inverse $2$-riffle law
\begin{equation}
\mathbb P(\pi)=2^{-n}\binom{n+1-d(\pi)}{n},                \tag{5}
\end{equation}
where $d(\pi)$ is the number of descents;
\item Ewens measure with $\theta=1.7$,
$\mathbb P(\pi)\propto\theta^{K(\pi)}$;
\item the general cycle weight
$\mathbb P(\pi)\propto1.2^{C_1}0.7^{C_2}1.6^{C_3}$.
\end{enumerate}
For $S_3,S_4,S_5$ and every partition through total degree four, the direct
route averages characters computed from the explicit matrices using Newton's
identities.  The cycle route applies (4) to the cycle structure.  Thus the two
routes share the probability law but not the spectral computation.

The normalizing generating-function formulas follow immediately from the
exponential formula for labeled cycles.  For example,
\begin{equation}
\Sfun_{n,\theta}^{\mathrm{Ewens}}=
\frac{[z^n]\exp\!\left(\theta\sum_{d\geq1}Q_dz^d/d\right)}
{[z^n]\exp\!\left(\theta\sum_{d\geq1}z^d/d\right)}.        \tag{6}
\end{equation}
Equation (6) records the corrected typography: the normalization is
$(\theta)_n$, and the summation index is $d\geq1$.

\section{Fourier proof and an independent \texorpdfstring{$S_3$}{S3} test}
Decompose
\begin{equation}
W_\tau=\bigotimes_a\operatorname{Sym}^{\tau_a}V
\cong\bigoplus_{\lambda\vdash n}M_{\tau,\lambda}\otimes V_\lambda,
\qquad m_{\tau,\lambda}=\dim M_{\tau,\lambda}.            \tag{7}
\end{equation}
Linearity of trace gives, for an arbitrary measure,
\begin{equation}
[\mathbf t^\tau]\Sfun_{\nu,n}
=\sum_{\lambda\vdash n}m_{\tau,\lambda}
\Tr\widehat\nu(\lambda),
\qquad
\widehat\nu(\lambda)=\sum_\pi\nu(\pi)\rho_\lambda(\pi). \tag{8}
\end{equation}
For a convolution walk, $\widehat{\kappa^{*r}}(\lambda)
=\widehat\kappa(\lambda)^r$.  If $\kappa$ is central, Schur's lemma makes
this operator scalar.  In the lazy transposition case,
\begin{equation}
\widehat\kappa(\lambda)=\beta_\lambda I,
\qquad
\beta_\lambda=a+(1-a)\frac{\chi^\lambda((12))}{d_\lambda}. \tag{9}
\end{equation}
For adjacent transpositions the full matrix power must be retained.

The $S_3$ test constructs the trivial, sign, and two-dimensional standard
matrices explicitly.  Multiplicities in (7) are obtained independently by
the uniform character inner product, then (8) is evaluated for an arbitrary
law, lazy central and adjacent walks, and the riffle law.  For laziness
$a=0.35$, the central scalars on $(\mathbf1,\mathrm{sgn},\mathrm{std})$ are
\begin{equation}
(1,-0.3,0.35).                                               \tag{10}
\end{equation}
The adjacent operator on the standard representation instead has two distinct
eigenvalues $0.025$ and $0.675$, directly exposing the claimed matrix-valued
behavior.

\section{Results}
There are 12 partitions through total degree four.  The cycle calculation
therefore contributes $3\cdot6\cdot12=216$ comparisons.  The four $S_3$
Fourier laws contribute 48 more, and the three irreducible operator checks
bring the total to 267.
\begin{center}
\begin{tabularx}{\linewidth}{@{}lXrrc@{}}
\toprule
Case & coefficient & matrix/direct & proposed formula & result\\
\midrule
$S_3$, inverse $2$-riffle & $(3)$ & $5.750000$ & $5.750000$ & \PASS\\
$S_4$, Ewens $\theta=1.7$ & $(2,1)$ & $7.429557$ & $7.429557$ & \PASS\\
$S_5$, transposition walk & $(2,2)$ & $39.683180$ & $39.683180$ & \PASS\\
$S_3$, arbitrary Fourier law & $(2,1)$ & $1.904762$ & $1.904762$ & \PASS\\
$S_3$, adjacent walk Fourier law & $(2,2)$ & $9.671793$ & $9.671793$ & \PASS\\
\bottomrule
\end{tabularx}
\end{center}

\section{Warnings, corrections, and scope}
Without laziness, a transposition walk alternates parity and does not converge
to uniform measure on $S_n$.  A walk on odd-length cycles stays in $A_n$.
These are support obstructions, not failures of (3) or (8).

The supplied continuous-time class-walk expressions contained a stray comma
in the exponential; the intended factor is
$\exp(t(r_\lambda-1))$.  Continuous-time graph interchange processes were not
numerically tested here, but their operator formula follows from replacing
matrix powers in (8) by the matrix exponential of the generator.

Run
\begin{verbatim}
.venv/bin/python -m \
  lambda_distributions.proofs.nonuniform_lambda_distributions.symmetric_groups.verification
\end{verbatim}
or open \path{marimo_notebooks/nonuniform_symmetric.py} in marimo.
\end{document}
